% Generated draft: requires mathematical and source review.
\documentclass[11pt]{article}
\usepackage[margin=25mm]{geometry}
\usepackage{fontspec}
\setmainfont{Latin Modern Roman}
\usepackage{amsmath,amssymb,mathtools,unicode-math}
\setmathfont{Latin Modern Math}
\usepackage{enumitem,needspace,xurl,hyperref}
\setlength{\parindent}{0pt}
\setlength{\parskip}{5pt}
\setlength{\emergencystretch}{3em}
\newcommand{\N}{\mathbb{N}}
\newcommand{\Z}{\mathbb{Z}}
\newcommand{\Q}{\mathbb{Q}}
\newcommand{\R}{\mathbb{R}}
\newcommand{\C}{\mathbb{C}}
\newcommand{\F}{\mathbb{F}}
\newcommand{\A}{\mathbb{A}}
\newcommand{\PP}{\mathbb{P}}
\newcommand{\E}{\mathbb{E}}
\newcommand{\eps}{\varepsilon}
\newcommand{\dd}{\,\mathrm{d}}
\begin{document}
\section*{UnsolvedMath ID: 30003907}
\textbf{Source problem code:} catalogue entry
\subsection{Definitions and mathematical statement}
A \(q\)-valued function takes values in unordered \(q\)-tuples, with multiplicities allowed. For a target set \(X\), write \(\mathcal A_q(X)=X^q/S_q\), where \(S_q\) acts by permuting the entries. The supplied record does not specify the domain, target geometry, admissible values of \(q\), or the Sobolev and comparison classes used to define Dirichlet-energy minimization. It also does not define a singular point or the centering, amplitude normalization, and convergence topology used in taking blow-ups.

The sequence-independence question can nevertheless be stated precisely once those choices are fixed. For an admissible minimizer \(u\) and a singular point \(Y\), let \(B_{Y,\rho}u\) denote the prescribed blow-up at scale \(\rho>0\), and let \(\tau\) denote the prescribed convergence topology. Define the set of tangent functions by
\[
\mathcal T_\tau(u,Y)=\left\{T:\ \text{there exists }\rho_j\downarrow0\text{ such that }B_{Y,\rho_j}u\longrightarrow T\text{ in }\tau\right\}.
\]
The supplied problem asks whether, for every such \(u\) and every singular point \(Y\), all tangent functions agree. Explicitly, for any two sequences \(\rho_j\downarrow0\) and \(\sigma_j\downarrow0\), is it true that
\[
B_{Y,\rho_j}u\longrightarrow T_1\text{ in }\tau,
\qquad
B_{Y,\sigma_j}u\longrightarrow T_2\text{ in }\tau
\quad\Longrightarrow\quad T_1=T_2?
\]
Equivalently, does \(\lvert\mathcal T_\tau(u,Y)\rvert\leq1\) always hold? If existence of a tangent is part of the intended background hypotheses, this becomes \(\lvert\mathcal T_\tau(u,Y)\rvert=1\). Uniqueness of subsequential limits alone does not assert convergence along every scale sequence without an appropriate compactness hypothesis. The missing blow-up conventions and analytic hypotheses must be recovered before this template specifies a complete mathematical conjecture.

The supplied dated assessment states that uniqueness of cylindrical tangent functions is proved at \(\mathcal H^{n-2}\)-almost every singular point, while uniqueness at every point or for every blow-up sequence remains open. Here \(\mathcal H^{n-2}\) denotes Hausdorff measure in the source's dimension notation. The record does not specify \(n\), define cylindrical tangents, or supply the precise hypotheses and quantifiers of that reported result; this assessment is retained as supplied curation, not asserted as an independently verified theorem.
\subsection{Short English statement}
For every Dirichlet-energy-minimizing \(q\)-valued function and every singular point, must all tangent functions obtained along scales tending to zero coincide? The domain, target, singularity definition, blow-up normalization, and convergence topology require clarification. The supplied record labels the problem partially solved, with a dated claim concerning cylindrical tangents at almost every singular point.
\subsection{Sources}
\begin{itemize}
\item Local UnsolvedMath registry, \texttt{data/unsolvedmath/problems.json}, supplied record with permanent ID \(30003907\), title ``Uniqueness of Tangents for Multivalued Energy Minimizers,'' exact source code \texttt{catalogue entry}, in the Oberwolfach Reports Open Problems collection. The record gives proposed year \(2018\) and identifies no individual proposer.
\item Supplied report citation: \emph{Calculus of Variations (2019)}, Oberwolfach Reports, DOI \texttt{10.4171/owr/2018/35}, pp.\ 2077--2156. \url{https://doi.org/10.4171/owr/2018/35}. The background attributes the problem to workshop open-problem discussions. The citation's year \(2019\) is retained alongside the DOI's \(2018\) component without resolving their relationship.
\item Supplied literature source: Krummel--Wickramasekera 2017. \url{https://arxiv.org/abs/1711.06222}. No paper title, theorem number, or exact hypotheses are supplied.
\item Supplied literature source: Oberwolfach Report 35. \url{https://doi.org/10.4171/owr/2018/35}.
\item Literature assessment date: \texttt{2026-08-21}. The supplied dated curation labels the problem \texttt{partially\_solved} and reports uniqueness of cylindrical tangent functions at \(\mathcal H^{n-2}\)-almost every singular point, with the universal uniqueness question remaining open. This is supplied dated curation, not independently verified; the source itself requests independent verification of time-sensitive status and cited claims.
\item The registry labels the statement \texttt{corrected\_verified} and explains that mathematical content was preserved while notation, typography, and sentence structure were repaired for publication. This supplied editorial label does not establish mathematical verification of the present draft.
\end{itemize}
\end{document}
